\section{D. Irga B. Naufal Fakhri}
\subsection{Soal 1}
Buatlah fungsi untuk membuka file csv dengan lib csv mode list
\lstinputlisting[firstline=9, lastline=20]{src/4/1174066/Praktek/lib_1174066_csv.py}

\subsection{Soal 2}
Buatlah fungsi untuk membuka file csv dengan lib csv mode dictionary
\lstinputlisting[firstline=24, lastline=34]{src/4/1174066/Praktek/lib_1174066_csv.py}

\subsection{Soal 3}
Buatlah fungsi untuk membuka file csv dengan lib pandas mode list
\lstinputlisting[firstline=9, lastline=11]{src/4/1174066/Praktek/lib_1174066_pandas.py}

\subsection{Soal 4}
Buatlah fungsi untuk membuka file csv dengan lib pandas mode dictionary
\lstinputlisting[firstline=15, lastline=19]{src/4/1174066/Praktek/lib_1174066_pandas.py}

\subsection{Soal 5}
Buat fungsi baru untuk mengubah format tanggal menjadi standar dataframe
\lstinputlisting[firstline=22, lastline=24]{src/4/1174066/Praktek/lib_1174066_pandas.py}

\subsection{Soal 6}
Buat fungsi baru  untuk mengubah index kolom
\lstinputlisting[firstline=28, lastline=30]{src/4/1174066/Praktek/lib_1174066_pandas.py}

\subsection{Soal 7}
Buat fungsi baru untuk mengubah atribut atau nama kolom
\lstinputlisting[firstline=35, lastline=37]{src/4/1174066/Praktek/lib_1174066_pandas.py}

\subsection{Soal 8}
Buat program main yang menggunakan library NPM csv yang membuat dan membaca file csv
\lstinputlisting[caption=lib\textunderscore 1174066\textunderscore csv.py, firstline=38, lastline=43]{src/4/1174066/Praktek/lib_1174066_csv.py}

\lstinputlisting[caption=main.py, firstline=8, lastline=17]{src/4/1174066/Praktek/main.py}

\subsection{Soal 9}
Buat program main2.py yang menggunakan library NPM pandas.py yang membuat dan membaca file csv
\lstinputlisting[caption=main2.py, firstline=42, lastline=48]{src/4/1174066/Praktek/lib_1174066_pandas.py}
\lstinputlisting[caption=main2.py, firstline=8, lastline=17]{src/4/1174066/Praktek/main2.py}

\subsection{Keterampilan Penanganan Error}
Pada praktikum saat ini saya tidak mendapatkan error

%%%%%%%%%%%%%%%%%%%%%%%%%%%%%%%%%%%%%%%%%%%%%%%%%%%%%%%%%%%%%%%%%%%%%%%%%%%%%%%%%%%%%%%%%%%%%%%%%%%%%%%%%%%%%%%%%%%%%%%%%%
\section{Fanny Shafira Damayanti | 1174069}
\subsection{Keterampilan Pemrograman}
\begin{enumerate}
	\item NO 1 

	\lstinputlisting[firstline=10, lastline=15]{src/4/1174069/Praktek/1174069_csv.py}

	\item NO 2

	\lstinputlisting[firstline=17, lastline=22]{src/4/1174069/Praktek/1174069_csv.py}

	\item NO 3 

	\lstinputlisting[firstline=10, lastline=13]{src/4/1174069/Praktek/1174069_pandas.py}

	\item NO 4 

	\lstinputlisting[firstline=10, lastline=13]{src/4/1174069/Praktek/1174069_pandas.py}

	\item NO 5 

	\lstinputlisting[firstline=15, lastline=19]{src/4/1174069/Praktek/1174069_pandas.py}

	\item NO 6 

	\lstinputlisting[firstline=21, lastline=24]{src/4/1174069/Praktek/1174069_pandas.py}

	\item NO 7 

	\lstinputlisting[firstline=26, lastline=30]{src/4/1174069/Praktek/1174069_pandas.py}

	\item NO 8

	\lstinputlisting[firstline=8, lastline=13]{src/4/1174069/Praktek/main.py}

	\item NO 9 

	\lstinputlisting[firstline=8, lastline=13]{src/4/1174069/Praktek/main2.py}

\end{enumerate}

\subsection{Penanganan Error}
\begin{enumerate}
	\item Peringatan error yang terdapat pada praktikum chapter 4 ini yaitu :

	\begin{itemize}
		\item Syntax Errors
		Syntax Errors terjadi ketika ada kesalahan dalam meuliskan kode. Solusinya adalah memperbaiki penulisan kode yang salah.

		\item Name Error
		NameError terjadi ketika salah mengetikan kode local name yang tidak terdefinisi. Solusinya adalah menuliskan kode dengan benar agar function nya dapat terpanggil. 

		\item Type Error
		TypeError terjadi pada saat eksekusi terhadapt fungsi dengan tipe objek tidak sesuai. Solusinya mengkonversi variablenya harus sesuai dengan tipe datanya.
	\end{itemize}

	Contoh Penggunaan TryExcept
	\lstinputlisting[firstline=55, lastline=67]{src/4/1174069/Praktek/1174069.py}
\end{enumerate}

%%%%%%%%%%%%%%%%%%%%%%%%%%%%%%%%%%%%%%%%%%%%%%%%%%%%%%%%%%%%%%%%%%%%%%%%%%%%%%%%%%%%%%%%%%%%%%%%%%%%%%%%%%%%%%%%%%%%
\section{Aulyardha Anindita | 1174054}
\subsection{Keterampilan Pemrograman}
\begin{enumerate}

\item Jawaban Soal No. 1
\lstinputlisting[firstline=10, lastline=15]{src/4/1174054/Praktek/1174054csv.py}

\item Jawaban Soal No. 2
\lstinputlisting[firstline=17, lastline=22]{src/4/1174054/Praktek/1174054csv.py}

\item Jawaban Soal No. 3
\lstinputlisting[firstline=10, lastline=13]{src/4/1174054/Praktek/1174054pandas.py}

\item Jawaban Soal No. 4
\lstinputlisting[firstline=10, lastline=13]{src/4/1174054/Praktek/1174054pandas.py}

\item Jawaban Soal No. 5
\lstinputlisting[firstline=15, lastline=19]{src/4/1174054/Praktek/1174054pandas.py}

\item Jawaban Soal No. 6
\lstinputlisting[firstline=21, lastline=24]{src/4/1174054/Praktek/1174054pandas.py}

\item Jawaban Soal No. 7
\lstinputlisting[firstline=26, lastline=30]{src/4/1174054/Praktek/1174054pandas.py}

\item Jawaban Soal No. 8
\lstinputlisting[firstline=8, lastline=13]{src/4/1174054/Praktek/main.py}

\item Jawaban Soal No. 9
\lstinputlisting[firstline=8, lastline=13]{src/4/1174054/Praktek/main2.py}

\subsection{Keterampilan Penanganan Error}

Peringatan error di praktek keempat ini, yaitu:
\begin{itemize}
\item Syntax Errors
Syntax Errors adalah keadaan dimana pada kode python mnengalami kesalahan dalam penulisan. Untuk mengatasinya yaitu dengan memperbaiki penulisan kode yang salah 

\item Name Error
NameError adalah suatu keadaan atau exception yang terjadi ketika kode melakukan eksekusi terhadap local name atau global name yang tidak terdefinisi. Untuk mengatasinya yaitu dengan memastikan variabel atau function yang dipanggil ada atau tidak salah ketik.

\item Type Error
TypeError adalah suatu keadaan atau exception yang akan terjadi apabila pada saat dilakukannya eksekusi terhadap suatu operasi atau fungsi dengan type object yang tidak sesuai. Untuk mengatasinya yaitu dengan mengkoversi varibelnya sesuai dengan tipe data yang akan digunakan.
\end{itemize}

Fungsi yang menggunakan try except
\lstinputlisting[firstline=55, lastline=67]{src/4/1174054/Praktek/1174054.py}

\end{enumerate}

%%%%%%%%%%%%%%%%%%%%%%%%%%%%%%%%%%%%%%%%%%%%%%%%%%%%%%%%%%%%%%%%%%%%%%

\section{Nurul Izza Hamka | 1174062 | Teori}
\section{Keterampilan Pemrograman}
\begin{enumerate}

\item NO1

\lstinputlisting[firstline=10, lastline=15]{src/4/1174062/Praktek/1174062csv.py}

\item NO2

\lstinputlisting[firstline=17, lastline=22]{src/4/1174062/Praktek/1174062csv.py}

\item NO3
	
\lstinputlisting[firstline=10, lastline=13]{src/4/1174062/Praktek/1174062Pandas.py}

\item NO4

\lstinputlisting[firstline=10, lastline=13]{src/4/1174062/Praktek/1174062Pandas.py}

\item NO5

\lstinputlisting[firstline=15, lastline=19]{src/4/1174062/Praktek/1174062Pandas.py}

\item NO6

\lstinputlisting[firstline=21, lastline=24]{src/4/1174062/Praktek/1174062Pandas.py}

\item NO7

\lstinputlisting[firstline=26, lastline=30]{src/4/1174062/Praktek/1174062Pandas.py}

\item NO8

\lstinputlisting[firstline=8, lastline=13]{src/4/1174062/Praktek/main.py}

\item NO9

\lstinputlisting[firstline=8, lastline=13]{src/4/1174062/Praktek/main2.py}

\section{Ketrampilan Penanganan Error}

\begin{itemize}
\item NameError
Error yang terjadi adalah NameError yang mana terjadi kesalahan pada saat melakukan eksekusi dan tidak dapat terdefinisi.
Penanganan yang dapat dilakukan adalah memastikan bahwa Variable dan Function yang akan kita panggil pastikan semua benar.

\item SyntaxError
Tipe Error ini adalah ada kesalahan pada penulisan, untuk itu pastikan penulisannya benar.

\item TypeError
Error ini terjadi pada saat eksekusi terhadap dan operasi dan fungsi masuk kedalam objek sedangkan typenya salah.

\item Menggunakan Try Except 

\lstinputlisting[firstline=55, lastline=67]{src/4/1174062/Praktek/1174062.py}
\end{itemize}



%%%%%%%%%%%%%%%%%%%%%%%%%%%%%%%%%%%%%%%%%%%%%%%%%%%%%%%%%%%%%%%%%%%%%%%%%%%%%%%%%%%%%%%%%%%%%%%%%%%%%%%%%%%%%%%%%%%%%%%%%%%%%%%%%%%%%%%%%%%%%%%%%%%%%%%%%%%%%%%%%%%%%%%%%%%%%%%%%%%%%%%%%%%%%%%%%%%
\section{Tia Nur Candida}
\begin{enumerate}
	\item Buatlah  fungsi  (file  terpisah/library  dengan  nama  NPMcsv.py)  untuk  membuka file csv dengan lib csv mode list.
	
	\lstinputlisting[caption = Fungsi untuk membuka file CSV dengan lib CSV mode list., firstline=10, lastline=15]{src/4/1174086/1174086csv.py}
	
	\item Buatlah  fungsi  (file  terpisah/library  dengan  nama  NPMcsv.py)  untuk  membuka file csv dengan lib csv mode dictionary.
	
	\lstinputlisting[caption =  Fungsi untuk membuka file CSV dengan lib CSV mode dictionary., firstline=17, lastline=22]{src/4/1174086/1174086csv.py}
	
	\item Buatlah fungsi (file terpisah/library dengan nama NPMpandas.py) untuk membuka file csv dengan lib pandas mode list.
	
	\lstinputlisting[caption =  Fungsi untuk membuka file CSV dengan lib Pandas mode list., firstline=10, lastline=13]{src/4/1174086/1174086pandas.py}
	
	\item Buatlah fungsi (file terpisah/library dengan nama NPMpandas.py) untuk membuka file csv dengan lib pandas mode dictionary.
	
	\lstinputlisting[caption =  Fungsi untuk membuka file CSV dengan lib Pandas mode dictionary., firstline=10, lastline=13]{src/4/1174086/1174086pandas.py}
	
	\item  Buat fungsi baru di NPMpandas.py untuk mengubah format tanggal menjadi standar dataframe.
	
	\lstinputlisting[caption =  Fungsi untuk mengubah format tanggal menjadi standar dataframe., firstline=15, lastline=19]{src/4/1174086/1174086pandas.py}
	
	\item Buat fungsi baru di NPMpandas.py untuk mengubah index kolom.
	
	\lstinputlisting[caption =  Fungsi untuk mengubah index kolom., firstline=21, lastline=24]{src/4/1174086/1174086pandas.py}
	
	\item Buat fungsi baru di NPMpandas.py untuk mengubah atribut atau nama kolom.
	
	\lstinputlisting[caption =  Fungsi untuk mengubah atribut atau nama kolom., firstline=26, lastline=30]{src/4/1174086/1174086pandas.py}
	
	\item Buat program main.py yang menggunakan library NPMcsv.py yang membuat dan membaca file csv.
	
	\lstinputlisting[caption =  Membuat dan mebaca file CSV menggunakan library 1174086pandas., firstline=8, lastline=13]{src/4/1174086/main.py}
	
	\item Buat program main2.py yang menggunakan library NPMpandas.py yang membuat dan membaca file csv.
	
	\lstinputlisting[caption = Membuat dan mmebaca file CSV menggunakan library 1174086pandas., firstline=8, lastline=13]{src/4/1174086/main2.py}
	
\end{enumerate}

\textbf{Kode Program}
\begin{figure}[H]
	\includegraphics[width=10cm]{figures/4/1174086/k1.png}
	\centering
\end{figure}
\begin{figure}[H]
	\includegraphics[width=10cm]{figures/4/1174086/k2.png}
	\centering
\end{figure}
\begin{figure}[H]
	\includegraphics[width=10cm]{figures/4/1174086/k3.png}
	\centering
\end{figure}
\begin{figure}[H]
	\includegraphics[width=10cm]{figures/4/1174086/k4.png}
	\centering
\end{figure}
\begin{figure}[H]
	\includegraphics[width=10cm]{figures/4/1174086/k5.png}
	\centering
\end{figure}

\textbf{Cek Plagiat}
\begin{figure}[H]
	\includegraphics[width=10cm]{figures/4/1174086/plagiatketrampilan.png}
	\centering
\end{figure}
\section{Penanganan Error}
	\begin{enumerate}
	\item Tuliskan  peringatan  error  yang  didapat  dari  mengerjakan  praktek  keempat  ini, dan  jelaskan  cara  penanganan  error  tersebut.   dan  Buatlah  satu  fungsi  yang menggunakan gunakan try except untuk menanggulangi error tersebut.
	\end{enumerate}
	Peringatan error di praktek keempat ini, yaitu:
	\begin{itemize}
		\item Syntax Errors
		Syntax Errors adalah suatu keadaan saat kode python mengalami kesalahan penulisan. Solusinya adalah memperbaiki penulisan kode yang salah.
		
		\item Name Error
		NameError adalah exception yang terjadi saat kode melakukan eksekusi terhadap local name atau global name yang tidak terdefinisi. Solusinya adalah memastikan variabel atau function yang dipanggil ada atau tidak salah ketik.
		
		\item Type Error
		TypeError adalah exception yang akan terjadi apabila pada saat dilakukannya eksekusi terhadap suatu operasi atau fungsi dengan type object yang tidak sesuai. Solusi dari error ini adalah mengkoversi varibelnya sesuai dengan tipe data yang akan digunakan.
	\end{itemize}
	
	Fungsi yang menggunakan try except
	\lstinputlisting[caption= Fungsi yang menggunakan try except .,firstline=55, lastline=67]{src/4/1174086/1174086.py}
\end{enumerate}

\textbf{Kode Program}
\begin{figure}[H]
	\includegraphics[width=10cm]{figures/4/1174086/p1.png}
	\centering
\end{figure}

\textbf{Cek Plagiat}
\begin{figure}[H]
	\includegraphics[width=10cm]{figures/4/1174086/plagiatpenanganan.png}
	\centering
\end{figure}

%%%%%%%%%%%%%%%%%%%%%%%%%%%%%%%%%%%%%%%%%%%%%%%%%%%%%%%%%%%%%%%%%%%%%%%%%%%%%%%%%%%%%%%%%%%%%%%%%%%%%%%

%%%%%%%%%%%%%%%%%%%%%%%%%%%%%%%%%%%%%%%%%%%%%%%%%%%%%%%%%%%%%%%%%%%%%%%%

\section{Difa Al Fansha}
\subsection{Soal Nomor 1}
Buatlah fungsi untuk membuka file csv dengan lib csv mode list

\lstinputlisting[firstline=10, lastline=17]{src/4/1174076/Praktek/coba.csv}
\begin{figure}[!htbp]
  \centering
  \includegraphics[height=3cm]{figures/4/1174076/Praktek/coba.png}
  \caption{Tampilan Kode Coba.scv}
\end{figure}

\lstinputlisting[firstline=10, lastline=17]{src/4/1174076/Praktek/1174076_csv.py}
\begin{figure}[!htbp]
  \centering
  \includegraphics[height=3cm]{figures/4/1174076/Praktek/1174076_csv.png}
  \caption{Output csv mode directoy, untuk Nomor 1,2,dan 8}
\end{figure}


\subsection{Soal Nomor 2}
Buatlah fungsi untuk membuka file csv dengan lib csv mode dictionary
\lstinputlisting[firstline=21, lastline=27]{src/4/1174076/Praktek/1174076_csv.py}

\subsection{Soal Nomor 3}
Buatlah fungsi untuk membuka file csv dengan lib pandas mode list
\lstinputlisting[firstline=11, lastline=15]{src/4/1174076/Praktek/1174076_pandas.py}
\begin{figure}[!htbp]
  \centering
  \includegraphics[height=3cm]{figures/4/1174076/Praktek/1174076_pandas.png}
  \caption{Output pandas mode list Nomor 3,4,5,6,7 dan 9}
\end{figure}

\subsection{Soal Nomor 4}
Buatlah fungsi untuk membuka file csv dengan lib pandas mode dictionary
\lstinputlisting[firstline=18, lastline=23]{src/4/1174076/Praktek/1174076_pandas.py}

\subsection{Soal Nomor 5}
Buat fungsi baru untuk mengubah format tanggal menjadi standar dataframe
\lstinputlisting[firstline=26, lastline=30]{src/4/1174076/Praktek/1174076_pandas.py}

\subsection{Soal Nomor 6}
Buat fungsi baru untuk mengubah index kolom
\lstinputlisting[firstline=33, lastline=38]{src/4/1174076/Praktek/1174076_pandas.py}

\subsection{Soal Nomor 7}
Buat fungsi baru untuk mengubah atribut atau nama kolom
\lstinputlisting[firstline=41, lastline=46]{src/4/1174076/Praktek/1174076_pandas.py}

\subsection{Soal Nomor 8}
Buat program main yang menggunakan library NPM csv yang membuat dan membaca file csv
\lstinputlisting[firstline=8, lastline=14]{src/4/1174076/Praktek/main.py}
\begin{figure}[!htbp]
  \centering
  \includegraphics[height=3cm]{figures/4/1174076/Praktek/main.png}
  \caption{Output main.py}
\end{figure}

\lstinputlisting[firstline=30, lastline=39]{src/4/1174076/Praktek/1174076_csv.py}

\lstinputlisting[firstline=1, lastline=5]{src/4/1174076/Praktek/buatfile.csv}
\begin{figure}[!htbp]
  \centering
  \includegraphics[height=3cm]{figures/4/1174076/Praktek/buatfile.png}
  \caption{Tampilan Koe buatfile.csv}
\end{figure}


\subsection{Soal Nomor 9}
Buat program main2.py yang menggunakan library NPM pandas.py yang membuat dan membaca file csv

\lstinputlisting[firstline=8, lastline=14]{src/4/1174076/Praktek/main2.py}
\begin{figure}[!htbp]
  \centering
  \includegraphics[height=3cm]{figures/4/1174076/Praktek/main2.png}
  \caption{Output pandas mode list Nomor 3,4,5,6,7 dan 9}
\end{figure}

\lstinputlisting[firstline=30, lastline=39]{src/4/1174076/Praktek/1174076_csv.py}

%%%%%%%%%%%%%%%%%%%%%%%%%%%%%%%%%%%%%%%%%%%%%%%%%%%%%%%%%%%%%%%%%%%%%%%%%%%%%%%%%%%%%%%%%%%%%%%%
